\documentclass[12pt]{article}
\usepackage{amsmath,amsthm,amssymb}
\usepackage{geometry}
\usepackage{setspace}
\usepackage{hyperref}

\geometry{margin=1in}
\onehalfspacing

\title{\textbf{The Field Guide to Noiselords and Yarncrawlers}\\
\large Practical Identification, Diagnosis, and Termination}
\author{Flyxion}
\date{2025}

\begin{document}
\maketitle

\section*{Purpose}

This document is a field guide. It is written for operators, researchers, designers, and observers who must interact with complex semantic systems under real constraints. It assumes the correctness of Fixed-Point Causality and provides operational criteria for recognizing, diagnosing, and intervening in systems that exhibit runaway interpretation.

This guide does not argue. It instructs.

\section{The Objects of Study}

All semantic systems encountered in the field fall into one of three regimes. These regimes are not moral categories. They are dynamical phases.

A \emph{Yarncrawler} is a system capable of lawful semantic traversal. It guarantees reachability across its semantic space while preserving structural constraints. A Yarncrawler can explore without inflating meaning.

A \emph{Noiselord} is a Yarncrawler that has lost its stopping rule. It remains surjective but is no longer governed by invariance. It cannot distinguish completion from continuation.

A \emph{Resolved System} is a Yarncrawler whose traversal is constrained by Fixed-Point Causality. It halts when further interpretation produces no change.

The field guide concerns the transitions between these regimes.

\section{Recognition: How to Identify a Noiselord}

A Noiselord can be identified without inspecting internal representations.

If a system continues to produce new semantic states while all observable invariants remain unchanged, the system is in Noiselord mode.

If a system treats silence, ambiguity, or symmetry as prompts for further interpretation rather than as termination signals, the system is in Noiselord mode.

If a system cannot report what would cause it to stop interpreting, the system is already stopped in the only sense that matters.

The defining diagnostic is simple: repeated application of the system’s transformation produces difference without distinction.

\section{Recognition: How to Identify a Yarncrawler}

A Yarncrawler is identifiable by its capacity for traversal without inflation.

A Yarncrawler can map between distant semantic regions without inventing new dimensions. It reveals latent structure rather than creating spurious structure.

Crucially, a Yarncrawler can fail meaningfully. It can reach states where no admissible traversal exists without violating constraints.

A Yarncrawler without Fixed-Point Causality is not defective. It is incomplete.

\section{Failure Modes}

Noiselord failure does not present as error. It presents as fluency.

The system remains coherent locally while drifting globally. Interpretations become increasingly confident while becoming increasingly untethered from invariance.

In human cognition, this appears as rumination, obsession, or confabulation. In institutions, it appears as bureaucratic expansion. In AI systems, it appears as hallucination.

In all cases, the signature is identical: interpretation continues after meaning has converged.

\section{The Termination Test}

The termination test is invariant across domains.

Given a system state $\Psi$, apply the system’s admissible transformation $F$.

If $F(\Psi)$ is observationally indistinguishable from $\Psi$ under the system’s own resolution, then the system must halt.

If the system cannot detect this condition, it is not intelligent in the operational sense. It is merely active.

\section{Intervention: How to Stop a Noiselord}

A Noiselord cannot be stopped by restriction of inputs. It must be stopped internally.

The first intervention is invariance detection. The system must be able to test whether its own transformations alter its state.

The second intervention is cost assignment. Each transformation must incur a measurable expenditure of internal resources.

The third intervention is refusal. When cost increases without invariant change, the system must irreversibly discard the traversal.

Stopping is not suppression. It is commitment.

\section{Worldhood as a Diagnostic}

Worldhood is the observable consequence of refusal.

A system has worldhood if there exist interpretations it will not pursue, even though it could.

A system without refusal inhabits no world. It inhabits a phase space.

The presence of worldhood is the strongest empirical indicator that Fixed-Point Causality is functioning.

\section{Operational Summary}

If a system cannot stop, it is not broken. It is unfinished.

If a system stops arbitrarily, it is brittle.

If a system stops when further transformation produces no difference, it is aligned with reality.

\section*{Final Note}

This field guide is complete when it is no longer needed.

Once you can see Noiselords in operation, they cannot be unseen. Once you recognize Yarncrawlers, you stop blaming traversal for drift.

Stopping is not the opposite of intelligence. It is its signature.

\end{document}

